\section{Conclusion}

HyperPod turns three mismatched consumer laptops into a single local runtime for
Meta open-weight models, and keeps serving when one of them stops.

The measurements in this report establish four things on real hardware. A
Llama~3.2~3B model was placed across a Linux RTX~3050, a Windows RTX~4050 and an
M4 MacBook~Air, with the per-device byte counts read back from
\texttt{llama.cpp} rather than assumed. Stopping the Mac accelerator mid-request
caused the inference process to abort, after which HyperPod excluded the
unreachable endpoint, rebuilt a still-distributed pipeline across the two
survivors in \SI{33}{\second}, and returned the same HTTP request without
resubmission; under the same conditions unmodified \texttt{llama.cpp} failed the
request and did not restart. A Llama~3.1~8B model that does not fit in the
controller's \SI{4096}{MiB} GPU ran with \SI{3330}{MiB} of its weights on the
remote Windows GPU. And distributing over Wi-Fi cost roughly a factor of six in
decode throughput relative to the single local GPU.

The last of those is as important as the first three. This design trades
throughput for capacity and for recovery. On a network like ours the trade only
makes sense when the model does not fit on one machine, or when an interrupted
request matters more than its latency.

What we would build next follows from Section~\ref{sec:limitations}. Caching
model weights on each worker's disk ahead of time would remove the long first
load, which is currently the worst part of the experience. Adding measured
compute and link bandwidth to the scheduler would let it place work on the
machine that can actually do it fastest, rather than the one with the most free
memory. Preserving the KV cache across a rebuild would turn restart-and-retry
into genuine continuation. Controller election would remove the last single
point of failure. Authentication would make the system safe outside a trusted
network.

Every figure quoted here is reproducible from the committed logs under
\code{test/} in the repository, and the two physical failure experiments are
recorded in \code{test/phase6-three-node-recovery.log} and
\code{test/phase7-native-benchmark.log}.

\vspace{1em}
\noindent\textbf{Built with Llama.} Llama 3.1 and Llama 3.2 are licensed by Meta
under their respective Community License Agreements. \texttt{llama.cpp} is used
under the MIT License. Full notices are in \code{THIRD\_PARTY\_NOTICES.md}.
